\documentclass[12pt]{article}

\usepackage{bbm}
\usepackage{fancyhdr}
\usepackage{titling}
\usepackage{enumitem}
\usepackage{amsmath}
\usepackage{dsfont}
\usepackage{amsfonts}
\usepackage{amssymb}
\usepackage{amsthm}
\usepackage{geometry}
\usepackage{tikz}
\usepackage{algorithm}
\usepackage{algpseudocode}
\geometry{verbose,tmargin=1in,bmargin=1in,lmargin=1in,rmargin=1in}
\renewcommand{\vec}[1]{\mathbf{#1}}
\newcommand{\dec}{\text{dec}}
\newcommand{\divi}{\mathbin{\textnormal{div}}}
\title{Homework 1}
\author{Satvik Prasad}
\date{}
\newtheorem{lemma}{Lemma}[section]
\fancypagestyle{plain}{%
  \fancyhf{}
  \fancyhead[L]{\theauthor}
  \fancyhead[R]{\thedate}
}
\pagestyle{plain}
\makeatletter
\def\@maketitle{%
  \newpage
  \null
  \vskip 1em%
  \begin{center}%
    \let \footnote \thanks
    {\LARGE \@title \par}%
    \vskip 1em%
  \end{center}%
  \par
\vskip 1em}
\makeatother
\usepackage{cmbright}

\newcommand{\mat}[1]{\mathbf{#1}}

% ---------- Probability theory ----------
% Core operators (optional argument-free; use as \E{X}, \Var{X}, \Cov{X}{Y})
\DeclareMathOperator{\EOp}{\mathbb{E}}
\DeclareMathOperator{\VarOp}{Var}
\DeclareMathOperator{\CovOp}{Cov}
\DeclareMathOperator{\CorrOp}{Corr}

\renewcommand{\P}[1]{\mathbb{P}\left(#1\right)}        % \P{A}
\newcommand{\E}[1]{\EOp\left[#1\right]}                % \E{X}
\newcommand{\Var}[1]{\VarOp\left(#1\right)}            % \Var{X}
\newcommand{\Cov}[2]{\CovOp\left(#1,\,#2\right)}       % \Cov{X}{Y}
\newcommand{\Corr}[2]{\CorrOp\left(#1,\,#2\right)}     % \Corr{X}{Y}
\newcommand{\SD}[1]{\operatorname{SD}\left(#1\right)}  % \SD{X}

% Subscripted versions, e.g. \Esub{\theta}{X}, \Psub{\theta}{A}
\newcommand{\Esub}[2]{\EOp_{#1}\left[#2\right]}
\newcommand{\Psub}[2]{\mathbb{P}_{#1}\left(#2\right)}

% Conditioning, indicators, independence
\newcommand{\given}{\,\middle|\,}                      % \E{X \given Y}
\newcommand{\ind}[1]{\mathbbm{1}\left\{#1\right\}}     % \ind{X > 0}
\newcommand{\indep}{\perp\!\!\!\perp}                  % X \indep Y
\newcommand{\iid}{\overset{\text{iid}}{\sim}}          % X_i \iid \Normal(0,1)

% Convergence
\newcommand{\convp}{\xrightarrow{\;p\;}}               % in probability
\newcommand{\convd}{\xrightarrow{\;d\;}}               % in distribution
\newcommand{\convas}{\xrightarrow{\text{a.s.}}}        % almost surely
\newcommand{\convLp}[1]{\xrightarrow{L^{#1}}}          % in L^p

% Probability spaces
\newcommand{\F}{\mathcal{F}}                           % sigma-algebra
\newcommand{\B}{\mathcal{B}}                           % Borel sigma-algebra
\newcommand{\pspace}{(\Omega, \F, \mathbb{P})}

% Common distributions
\newcommand{\Normal}{\mathcal{N}}
\newcommand{\Bern}{\operatorname{Bernoulli}}
\newcommand{\Bin}{\operatorname{Binomial}}
\newcommand{\Pois}{\operatorname{Poisson}}
\newcommand{\Unif}{\operatorname{Uniform}}
\newcommand{\Exp}{\operatorname{Exp}}
\newcommand{\Geom}{\operatorname{Geometric}}
\newcommand{\Gam}{\operatorname{Gamma}}
\newcommand{\Beta}{\operatorname{Beta}}

% Number sets
\newcommand{\R}{\mathbb{R}}
\newcommand{\N}{\mathbb{N}}
\newcommand{\Z}{\mathbb{Z}}


\begin{document}
\begin{enumerate}
\item Let us have the probability space $(\Omega, \mathcal F, \mathbb P).$ Assume we have discrete random variables $X(\omega)$ and $Y(\omega).$ Consider the map $U : \Omega \to \mathbb R$ defined by
$$
U(\omega) = X(\omega) + Y(\omega).
$$

Then, consider the image $U(\Omega),$
$$
U(\Omega) = \{X(\omega) + Y(\omega) : \omega \in \Omega\}.
$$
Since $X(\Omega)$ and $Y(\Omega)$ are countable, $X(\Omega) \times Y(\Omega)$ is countable. Define the mapping $f: X(\Omega) \times Y(\Omega) \to X(\Omega) + Y(\Omega),$ such that
$$
f(a, b) = a + b,
$$
where $X(\Omega) + Y(\Omega) = \{x + y: x\in X(\Omega), y \in Y(\Omega)\}.$ For any $u \in X(\Omega) + Y(\Omega),$ by definition there exists $x \in X(\Omega), y \in Y(\Omega)$ such that $x + y = u,$ meaning $f(x, y) = u,$ and hence $f$ is surjective. Therefore, $X(\Omega) + Y(\Omega)$ is countable.
Finally, $U(\Omega) \subseteq X(\Omega) + Y(\Omega),$ since $X(\omega) + Y(\omega) \in X(\Omega) + Y(\Omega)$ for all $w \in \Omega.$ Thus, $U(\Omega)$ is countable.
Finally, consider the pre-image $U^{-1}(u)$ for any $u \in \mathbb R.$ This is the set such that
$$
    U^{-1}(u) = \{\omega : X(\omega) + Y(\omega) = u \}.
$$
Enumerating over values of $Y(\omega),$ we have
\begin{align*}
U^{-1}(u) &= \bigcup_{y \in Y(\Omega)} \{\omega : X(\omega) = u - y\} \\
&= \bigcup_{y \in Y(\Omega)} X^{-1}(u - y).
\end{align*}
Since $\mathcal F$ is closed under countable union, and $X^{-1}(u - y) \in \mathcal F,$ then $U^{-1}(u) \in \mathcal F.$
Therefore, since $U(\Omega)$ is countable and the preimage $U^{-1}(u)$ is measurable for all $u \in \mathbb R,$ $U$ is a discrete random variable.
\item \begin{enumerate}
    \item Assume we have some discrete random variable $X \sim \text{Geometric}(p)$ for $p < 1.$ Then, $P(X = x) = p(1-p)^{x-1}$ for $x = \left\{1, 2, \ldots \right\}.$ Then, 
    $$
    \mathbb P(X > a + b \mid X > b) = \dfrac{\mathbb P(X > a + b)}{\mathbb P(X > b)}.
    $$
    Finding a general formula for $P(X > x),$ we have
    \begin{align*}
    P(X > x) &= \sum_{i = x+ 1}^{\infty} p(1-p)^{i - 1} \\
    &= p\sum_{i = x + 1}^{\infty} (1-p)^{i - 1} \\
    &= p \sum_{i = 0}^{\infty} (1-p)^{i + x} \\
    &= p (1 - p)^{x} \sum_{i = 0}^{\infty} (1-p)^{i} \\
    &= p(1 - p)^{x} \cdot \dfrac{1}{p} = (1-p)^x.
    \end{align*}
    Substituting this into our conditional probability, we get 
    $$
    \mathbb P(X > a + b \mid X > b) = \dfrac{(1-p)^{a + b}}{(1-p)^b} = (1- p)^a = \mathbb P(X > a).
    $$
    Therefore, we can conclude that if $X \sim \text{Geometric}(p),$ $X$ is memoryless.
    \item Assume we have discrete memoryless $X$ supported on $\{1, 2, \ldots\}.$ Then, we have 
    $$
    \mathbb P(X > a + b \mid X > b) = \mathbb P(X > a),
    $$
    which occurs if and only if 
    $$
    \mathbb P(X > a + b) = \mathbb P(X > a) \cdot \mathbb P(X > b).
    $$
    Letting $G(x) = \mathbb P(X > x),$ we get the requirement 
    $$
    G(a + b) = G(a) \cdot G(b).
    $$
    Letting $b = 0$, we know $G(0) = 1.$ Letting $b = 1,$ we get
    $$
    G(a + 1) = G(a) \cdot G(1),
    $$
    which expanding the recursion, yields
    $$
    G(x) = G(1)^{x} = \mathbb P(X > x).
    $$
    Let $G(1) = 1 -p$ for some $p \in \mathbb R.$ Then, $$
    \mathbb P(X > x) = (1-p)^x.
    $$
    By the definition of a probability measure, we need $0 \leq \mathbb P(X > x) \leq 1,$ meaning $0 \leq p \leq 1.$ Since
    $$
    \mathbb P(X>x - 1) = \sum_{i = x}^{\infty} \mathbb P(X = i) = \mathbb P(X = x) + \mathbb P(X > x),
    $$
    we can conclude
    $$
    \mathbb P(X = x) = (1-p)^{x-1} - (1-p)^x = (1-p)^{x-1} \cdot p,
    $$
    for some $0 \leq p \leq 1.$
    Finally, for $\mathbb P(\Omega) = 1,$ we need $\mathbb P(X > 0) = 1,$ meaning $(1 - p)^{0} = 1,$ which is only the case when $p < 1.$
    Thus, this forces $X \sim \text{Geometric}(p)$ for some $p < 1$ as required.
\end{enumerate}
    \item We look to find the pmf $p_X(x)$ that maximises 
    $$
    H(X) = -\sum_{s \in S} p_X(s) \log(p_X(s)).
    $$
    Using the properties of the logarithm, we get 
    $$
    H(X) = \sum_{s \in S} p_X(s) \log\left(\dfrac{1}{p_X(s)}\right).
    $$
    First, notice that the random variable $X$ described by the uniform pmf $p_X(s) = 1/k$ yields
    $$
    H(X) = \sum_{s \in S} \dfrac{1}{k} \log k = \log k.
    $$
    We prove that this is optimal. This is equivalent to showing
    $$
    f(X) = H(X) - \log k \leq 0.
    $$
    Manipulating further, we get 
    \begin{align*}
    f(X) &= H(X) - \sum_{s \in S} p_X(s) \log k \\
    &= \sum_{s \in S} p_X(s) \left[
    \log \left(\dfrac{1}{p_X(s)}\right) - \log k
    \right] \\
    &= \sum_{s \in S} p_X(s) \log \left(\dfrac{1}{kp_X(s)}\right).
    \end{align*}
    Since $\log x \leq x - 1$ for all $x \in (0, \infty),$ we can conclude
    $$
    f(X) \leq \sum_{s \in S} p_X(s) \left(\dfrac{1}{kp_X(s)} - 1\right) = 1 - \sum_{s \in S} p_X(s) = 0.
    $$
    Thus, $$
    H(x) \leq \log k,
    $$
    and the uniform distribution $p_X(s) = 1/k$ achieves this equality, meaning $p_X(s) = 1/k$ maximises the entropy.
    \item \begin{enumerate}
    \item Let the $N$ balls be represented by $b_1, \ldots, b_N. $Define our sample space $\Omega$ as 
    $$
    \Omega = \{
        (b_{i_1}, b_{i_2}, \ldots, b_{i_n}) : 1 \leq i_1 <i _2 < \ldots < i_n \leq N
    \},
    $$
    and let $b(b_i) = 1$ if $b_i$ is black, and $0$ otherwise.
    Then, we can define discrete random variable $B: \Omega \to \mathbb R$ such that
    $$
    B(\omega) = \sum_{b_i \in \omega} b(b_i).
    $$
    Therefore, the preimage $B^{-1}(t)$ consists of all $\omega$ such that $$
    \sum_{b_i \in \omega} b(b_i) = t.
    $$
    Thus, let us count 
    $$
    \left|\left\{\omega : \sum_{b_i \in \omega} b(b_i) = t\right\}\right|,
    $$
    where we need to choose $t$ out of the $b$ black balls and then the remaining $n - t$ out of the $w$ white balls. Thus, 
    $$
    \left|\left\{\omega : \sum_{b_i \in \omega} b(b_i) = t\right\}\right| ={b \choose t} \cdot {w \choose n - t}.
    $$
    Since these outcomes are disjoint, and each outcome has probability ${N \choose n}^{-1}$
    $$
    \mathbb P(B = t) = \dfrac{{b \choose t} {w \choose n - t}}{{N \choose n}}.
    $$
    \item Let $b = pN$ and $w = (1 - p)N.$ Then, in the limit as $N \to \infty,$ we have
    $$
    \lim_{N \to \infty} \mathbb P(B = t) = \lim_{N \to \infty} \dfrac{{Np \choose t} {N - Np \choose n - t}}{{N \choose n}}.
    $$
    Rearranging the inside of the limit yields
    $$
    {Np \choose t} \cdot t! \cdot \dfrac{1}{(Np)^t} \cdot {N - Np \choose n - t} \cdot (n - t)! \cdot \dfrac{1}{(N - Np)^{n - t}} \cdot \dfrac{(Np)^t \cdot (N - Np)^{n - t}}{t! (n-t)!} \cdot \dfrac{1}{{N \choose n}}.
    $$
    Via the given identity, the limit of this term as $N \to \infty$ is equivalent to 
    \begin{align*}
    \lim_{N \to \infty} \dfrac{(Np)^t (N - Np)^{n-t}}{t! (n - t)! {N \choose n}} &= {n \choose t} \lim_{N \to \infty} \dfrac{N^{n}p^t(1-p)^{n-t}}{n !{N \choose n}} \\
    &={n \choose t} \lim_{N \to \infty} p^t (1-p)^{n-t} \left(
    \dfrac{{N \choose n}}{\frac{1}{n!}N^{n}}
    \right)^{-1} \\
    &= {n \choose t } p^{t}(1 - p)^{n -t} = \mathbb P(\text{Binomial}(n, p) = t),
    \end{align*}
    and hence 
    $$
    \lim_{N \to \infty} \mathbb P(B = t) = \mathbb P(\text{Binomial}(n, p) = t)
    $$
    as required.
    \end{enumerate}
    \item \begin{enumerate}
        \item Label the four vertices $a, b, c, d$ in clockwise order, starting from the top left corner. Then, we have 
        \begin{align*}
        \mathbb P(a = b = c = d = -1 \mid c = -1)
        &= \dfrac{e^{4\beta}}{\sum_{(a,b,d) \in \{-1,1\}^3} e^{\beta(ab-b-d+ad)}} \\
        &= \dfrac{e^{4\beta}}{\sum_{(a,b,d) \in \{-1,1\}^3} e^{\beta(b(a-1)+d(a-1))}}.
        \end{align*}
        Looking at the denominator, 
        \begin{align*}
        \sum_{a,b,d = \pm 1}e^{\beta(b(a-1)+d(a-1))}
        &= \sum_{a,b,d = \pm 1} e^{\beta b(a-1)} e^{\beta d(a-1)} \\
        &= \sum_{a = \pm 1} \sum_{b = \pm 1} \sum_{d = \pm 1} e^{\beta b(a-1)} e^{\beta d(a-1)} \\
        &= \sum_{a = \pm 1} \left( \sum_{b = \pm 1} e^{\beta b(a-1)} \right) \left(\sum_{d = \pm 1} e^{\beta d(a-1)}\right) \\
        &= \sum_{a = \pm 1} 2\cosh (\beta (a - 1)) \cdot 2 \cosh (\beta(a - 1)) \\
        &= \sum_{a = \pm 1} 4 \cosh (\beta(a-1))^2 \\
        &= 4 \cosh(0)^2 + 4 \cosh(-2 \beta)^2 \\
        &= 4 + 4 \cosh (2 \beta)^2 \\
        &= 4(1 + \cosh^2(2 \beta)),
        \end{align*}
        yielding 
        \begin{align*}
        \mathbb P(a = b = c = d = - 1 \mid c= -1) &= \dfrac{e^{4\beta}}{4(1 + \cosh^2(2\beta))} \\
        &= \dfrac{e^{4\beta}}{2(3 + \cosh(4\beta))}.
        \end{align*}
        \item Similarly, we have
        $$
        \mathbb P(a = b = c = d = -1 \mid c = -1) = \dfrac{(1/2)^4}{(1/2)} = 1/8.
        $$
        \item From part (a), we have
        \begin{align*}
        \mathbb P(a = b = c= d = -1 \mid c = -1) &= \dfrac{e^{4\beta}}{6 + e^{4\beta} + e^{-4\beta}}.
        \end{align*}
        As $\beta \to 0,$ this fraction approaches $
        1/8,
        $
        whereas when $\beta \to \infty,$ we can write
        $$
        p = \dfrac{1}{6e^{-4\beta} + 1 + e^{-8\beta}} \to 1
        $$
        as $\beta \to \infty.$
        At high temperatures the distribution of colors are independently distributed and less likely to all be one uniform color, whereas at low temperature given that a single vertex is blue, the probability that the rest is blue approaches $1,$ meaning its more likely to be one uniform color (i.e., magnetized).
    \end{enumerate}

\end{enumerate}
\end{document}